% Part: methods
% Chapter: proofs
% Section: using-definitions

\documentclass[../../../include/open-logic-section]{subfiles}

\begin{document}

\olfileid{mth}{prf}{def}

\olsection{العمل بالتعريفات}

قد تقدم أن إحاطة الناظر بالتعريفات الداخلة في البرهان، وإجراؤها على
وجهها، شرط لا غنى عنه؛ وهذا أصل يستحق مزيد بيان. وإنما وُضعت
التعريفات لتجمع خواص الأشياء وعلاقاتها في ألفاظ موجزة. فاللفظ المختصر
الموضوع يسمى \emph{المعرَّف}، والمعنى المبسوط الذي يدل عليه يسمى
\emph{المُعرِّف}. وكثيرًا ما نحتاج في البرهان إلى رد المعرَّف إلى أصل
وضعه، لأن العمل إنما يقع على البنية المنطقية للمُعرِّف، أي على العبارة
المطولة التي قام اللفظ المعرَّف مقامها. فإذا فُك التعريف ظهر موضع
الاستدلال ولم يبق مستورًا بالاختصار.

ولنبيّن ذلك بمثال؛ نريد أن نثبت القضية الآتية:

\begin{prop}
لأي مجموعتين~$\OLId{latin-looped}{A}$ و$\OLId{latin-looped}{B}$، يكون~$\OLId{latin-looped}{A} \cup \OLId{latin-looped}{B} = \OLId{latin-looped}{B} \cup \OLId{latin-looped}{A}$.
\end{prop}

ولا يصح الشروع حتى نعلم ما تطابق مجموعتين، أي ما يدل عليه «$=$» في
هذه المعادلة إذا كان طرفاها مجموعتين. والمجموعتان متطابقتان متى
اشتركتا في !!{element}s بأعيانها؛ فالتعريف الذي ينبغي بسطه هو:

\begin{defn}
تكون المجموعتان~$\OLId{latin-looped}{A}$ و$\OLId{latin-looped}{B}$ \emph{متطابقتين}،~$\OLId{latin-looped}{A} = \OLId{latin-looped}{B}$، إذا وفقط إذا كان
كل !!{element} من~$\OLId{latin-looped}{A}$ !!a{element} من~$\OLId{latin-looped}{B}$، وصح العكس.
\end{defn}

وجيء في هذا التعريف بـ~$\OLId{latin-looped}{A}$ و$\OLId{latin-looped}{B}$ مثالين قائمين مقام أي مجموعتين. وأما
\emph{المعرَّف} فهو قولنا «$\OLId{latin-looped}{A} = \OLId{latin-looped}{B}$»؛ إذ يعيّن التعريف الشرط الذي يصدق
معه~$\OLId{latin-looped}{A} = \OLId{latin-looped}{B}$. وذلك الشرط، وهو أن «كل !!{element} من~$\OLId{latin-looped}{A}$ هو
!!a{element} من~$\OLId{latin-looped}{B}$، وصح العكس»، هو \emph{المُعرِّف}.\footnote{وفي
هذه الصورة خاصة موضع قد يوقع في الحيرة: فإذا كان~$\OLId{latin-looped}{A} = \OLId{latin-looped}{B}$، فالمذكوران
$\OLId{latin-looped}{A}$ و$\OLId{latin-looped}{B}$ مجموعة واحدة بعينها، وإن اختلف الحرف الدال عليها في الطرفين.
غير أن الطريق الذي عُيّنت به المجموعة في أحدهما قد يخالف الطريق الذي
عُيّنت به في الآخر، ومن هنا لم يكن التعريف تحصيلًا للحاصل.} فمؤدى
التعريف أن~$\OLId{latin-looped}{A} = \OLId{latin-looped}{B}$ يصدق متى تحقق الشرط، وألا يصدق إلا عند تحققه؛ وهذا
هو المراد بقولنا «إذا وفقط إذا».

وإذا أُجري التعريف على مسألة بعينها وجب تنزيل~$\OLId{latin-looped}{A}$ و$\OLId{latin-looped}{B}$ المذكورتين فيه
على طرفي تلك المسألة. ففي مثالنا لا يصدق~$\OLId{latin-looped}{A} \cup \OLId{latin-looped}{B} = \OLId{latin-looped}{B} \cup \OLId{latin-looped}{A}$ إلا
إذا كان كل~$\OLId{latin-ordinary}{z} \in \OLId{latin-looped}{A} \cup \OLId{latin-looped}{B}$ كائنًا أيضًا في~$\OLId{latin-looped}{B} \cup \OLId{latin-looped}{A}$، وصح العكس.
فـ~$\OLId{latin-looped}{A} \cup \OLId{latin-looped}{B}$ هنا بمنزلة~$\OLId{latin-looped}{A}$ في نص التعريف، و~$\OLId{latin-looped}{B} \cup \OLId{latin-looped}{A}$ بمنزلة~$\OLId{latin-looped}{B}$.
ولما كانت الحروف~$\OLId{latin-looped}{A}$ و$\OLId{latin-looped}{B}$ قد استُعملت في التعريف وفي القضية استعمالين
مختلفين، احتيج إلى التمييز بين المقامين. فلا يجوز، مثلًا، أن يُظن أن
إثبات كون كل !!{element} من~$\OLId{latin-looped}{A}$ !!a{element} من~$\OLId{latin-looped}{B}$ والعكس يثبت
القضية؛ فإن غاية ذلك إثبات~$\OLId{latin-looped}{A} = \OLId{latin-looped}{B}$، لا~$\OLId{latin-looped}{A} \cup \OLId{latin-looped}{B} = \OLId{latin-looped}{B} \cup \OLId{latin-looped}{A}$.
(ثم إن~$\OLId{latin-looped}{A}$ و$\OLId{latin-looped}{B}$ مجموعتان كيف اتفق، وقد تختلفان كل الاختلاف ما لم
يفرض شيء في~$\OLId{latin-looped}{A}$ و$\OLId{latin-looped}{B}$، فلا يفيد هذا الطريق أصلًا.)

وفي البرهان مفهوم مجموعاتي آخر هو الاتحاد، فلا بد من معرفة دلالة
$\cup$ قبل المضي. وقد يفضي بسط تعريف إلى تعريف آخر يجب بسطه أيضًا.
فـ~$\OLId{latin-looped}{A} \cup \OLId{latin-looped}{B}$، مثلًا، معرّفة بأنها
$\Setabs{\OLId{latin-ordinary}{z}}{\OLId{latin-ordinary}{z} \in \OLId{latin-looped}{A} \text{ أو } \OLId{latin-ordinary}{z} \in \OLId{latin-looped}{B}}$. فإذا أردنا إثبات
$\OLId{latin-ordinary}{x} \in \OLId{latin-looped}{A} \cup \OLId{latin-looped}{B}$، صار المطلوب بعد فك~$\cup$ إثبات
$\OLId{latin-ordinary}{x} \in \Setabs{\OLId{latin-ordinary}{z}}{\OLId{latin-ordinary}{z} \in \OLId{latin-looped}{A} \text{ أو } \OLId{latin-ordinary}{z} \in \OLId{latin-looped}{B}}$. ثم نرجع إلى أن
$\OLId{latin-ordinary}{x} \in \Setabs{\OLId{latin-ordinary}{z}}{\dots \OLId{latin-ordinary}{z}\dots}$ إذا وفقط إذا كانت~$\dots \OLId{latin-ordinary}{x}\dots$؛
فإذا بُسط أيضًا ترميز~$\Setabs{\OLId{latin-ordinary}{z}}{\dots \OLId{latin-ordinary}{z} \dots}$ آل المطلوب إلى
$\OLId{latin-ordinary}{x} \in \OLId{latin-looped}{A}$ أو~$\OLId{latin-ordinary}{x} \in \OLId{latin-looped}{B}$. وعلى ذلك فقولنا «كل !!{element} من
$\OLId{latin-looped}{A} \cup \OLId{latin-looped}{B}$ هو أيضًا !!a{element} من~$\OLId{latin-looped}{B} \cup \OLId{latin-looped}{A}$» مؤداه: «لكل~$\OLId{latin-ordinary}{x}$،
إذا كان~$\OLId{latin-ordinary}{x} \in \OLId{latin-looped}{A}$ أو~$\OLId{latin-ordinary}{x} \in \OLId{latin-looped}{B}$، فإن~$\OLId{latin-ordinary}{x} \in \OLId{latin-looped}{B}$ أو~$\OLId{latin-ordinary}{x} \in \OLId{latin-looped}{A}$».
وبسط جميع تعريفات القضية يردها إلى الصورة الآتية:

\begin{prop}
لكل مجموعتين~$\OLId{latin-looped}{A}$ و$\OLId{latin-looped}{B}$: (أ) لكل~$\OLId{latin-ordinary}{x}$، إذا كان~$\OLId{latin-ordinary}{x} \in \OLId{latin-looped}{A}$ أو
$\OLId{latin-ordinary}{x} \in \OLId{latin-looped}{B}$، كان~$\OLId{latin-ordinary}{x} \in \OLId{latin-looped}{B}$ أو~$\OLId{latin-ordinary}{x} \in \OLId{latin-looped}{A}$؛ و(ب) لكل~$\OLId{latin-ordinary}{x}$، إذا كان
$\OLId{latin-ordinary}{x} \in \OLId{latin-looped}{B}$ أو~$\OLId{latin-ordinary}{x} \in \OLId{latin-looped}{A}$، كان~$\OLId{latin-ordinary}{x} \in \OLId{latin-looped}{A}$ أو~$\OLId{latin-ordinary}{x} \in \OLId{latin-looped}{B}$.
\end{prop}

فبسط التعريفات، إذن، ركن من بناء البرهان لا تتم صناعته بدونه، وإن
دقّ أحيانًا مسلكه. إذ يلزم أن تُفصل متغيرات التعريف بعضها من بعض، ثم
يُعرف ما يقابل كلًّا منها من حدود الدعوى. ولا يتأتى ذلك حتى يُفهم
المطلوب وتُضبط دلالات حدوده جميعًا؛ وربما احتيج إلى بسط تعريفات
متتابعة. وفي البراهين السهلة، كالمسألة الآتية، يكاد الحل يخرج من
التعريف نفسه؛ وليس الأمر كذلك في كل موضع.

\begin{prob}
افرض أن المطلوب إثبات~$\OLId{latin-looped}{A} \cap \OLId{latin-looped}{B} \neq \emptyset$. ابسط التعريفات كلها؛
أي عبّر عن المطلوب من غير أن تذكر «$\cap$»، أو «$=$»، أو
«$\emptyset$».
\end{prob}

\end{document}
